\documentclass[12pt, letterpaper]{article}

%-------Packages---------
\usepackage{amssymb,amsfonts}
\usepackage{mathptmx}
\usepackage{amsthm}
\usepackage{graphicx}
\usepackage[all,arc]{xy}
\usepackage{enumerate}
\usepackage{bbm}
\usepackage{dsfont}
\usepackage{mathrsfs}
\usepackage{tikz-cd}
\usepackage{setspace}
\usepackage{import}
\usepackage[utf8]{inputenc}
\usepackage{hyperref}
\usepackage{tikz}
\usepackage{float}
\usepackage[dvipsnames]{xcolor}
\usepackage{fancyhdr}
\usepackage[nointegrals]{wasysym}

\usepackage[left=1in,top=1in,right=1in,bottom=1in]{geometry}

\usepackage{mathtools}
\usepackage{setspace}
\usepackage{cleveref}
\usepackage{multicol}
\usepackage{mathpazo}
\usepackage{tcolorbox}
\tcbuselibrary{theorems,skins,breakable}
\usepackage{xparse}


\usepackage{xcolor, ifthen, tcolorbox}
\tcbuselibrary{theorems,skins,breakable}
% Theorem System
% The following environments are provided:
%   New Problem:        \begin{problem} ...
%   New Question Header:   \qh (then followed by subproblems)
%   Subproblem:     \begin{subproblem} ...
%   Problem Proof:  \begin{problemproof} ...
%   Proof:          \\begin{pf}{color} ...
%   Remark:         \rmk (sentence), \rmkb (block)

% Define custom colors
\definecolor{thmscol}{HTML}{F4565B} %For Problems
\definecolor{lemscol}{HTML}{FE844C} %For Lemmes
\definecolor{prpscol}{HTML}{00A7EF} %For proposition
\definecolor{corscol}{HTML}{23DAFF} %For corrolaries
\definecolor{rmkscol}{HTML}{6DEEC5} %For remarks
\definecolor{defscol}{HTML}{18C991} %For definitions
\definecolor{exmscol}{HTML}{7DB800} %For examples
\definecolor{clmscol}{HTML}{165c58} %For claims

%Change between dark(black)/light(white) mode
\newcommand{\colormode}{white}
\ifthenelse{\equal{\colormode}{white}}
      {\newcommand{\TextColor}{black}}
      {\newcommand{\TextColor}{white}}
\newcommand{\pbcolormain}{thmscol!80!\colormode}
\newcommand{\pbcolorsecond}{thmscol!38!\colormode}


% ============================
% Problem Counter
% ============================
\newcounter{subproblemcounter}
\newcounter{problemcounter}

\newcommand{\q}{
    \setcounter{subproblemcounter}{0}%
    \stepcounter{problemcounter} % Increment problem counter
    \color{\TextColor}Problem \arabic{problemcounter}: % Display the problem counter
}

\newcommand{\stepsub}{
    \stepcounter{subproblemcounter}%
    \color{\TextColor}(\alph{subproblemcounter})
}
% ============================
% Problem
% ============================
\newtcolorbox{pb}[1]{
  enhanced,
  frame hidden,
  titlerule=0mm,
  toptitle=1mm,
  bottomtitle=1mm,
  fonttitle=\bfseries\large,
  coltitle=black,
  colbacktitle=\pbcolormain,
  colback=\pbcolorsecond,
  title={\q #1},
  coltext=\TextColor
}

\newtcolorbox{spb}[1]{
  enhanced,
  frame hidden,
  titlerule=0mm,
  toptitle=1mm,
  bottomtitle=1mm,
  fonttitle=\bfseries\large,
  coltitle=black,
  colbacktitle=\pbcolormain,
  colback=\pbcolorsecond,
  title={\stepsub #1},
  coltext=\TextColor,
  attach boxed title to top left={xshift=3mm,yshift=-2mm},
  boxed title style={
    colback=\pbcolormain,
    colframe=\pbcolormain,
  }
}

\newtcolorbox{pbh}[1]{
    enhanced,
    frame hidden,
    titlerule=0mm,
    toptitle=1mm,
    bottomtitle=1mm,
    fonttitle=\bfseries\large,
    coltitle=black,
    colbacktitle=\pbcolormain,
    colback=\pbcolorsecond,
    title={\q #1},
    boxrule=0pt,
    interior hidden,
    attach boxed title to top left={xshift=3mm,yshift=-2mm},
    boxed title style={
        colback=\pbcolormain,
        colframe=\pbcolormain,
    }
}

\newenvironment{problem}[1]{%
  \newpage
  \begin{pb}{#1}
    }{
  \end{pb}
}

\newenvironment{subproblem}[1]{%
  \begin{spb}{#1}
    }{
  \end{spb}
}

\newenvironment{problemproof}{%
    \begin{pf}{\pbcolormain}
        }{
    \end{pf}
}

\newcommand{\qh}[1][]{
    \newpage
    \begin{pbh}{#1}\end{pbh} % Display the problem counter
}

% ============================
% Claim
% ============================

\newtcbtheorem[number within=problemcounter]{claim}{Claim}
{
    enhanced,
    frame hidden,
    titlerule=0mm,
    toptitle=1mm,
    bottomtitle=1mm,
    fonttitle=\bfseries\large,
    coltitle=black,
    colbacktitle=clmscol!50!\colormode,
    colback=clmscol!20!\colormode,
}{clm}

\NewDocumentCommand{\clm}{m+m}{
    \begin{claim*}{#1}{}
        #2
    \end{claim*}
}

\newenvironment{clmpf}{
	{\noindent{\it \textbf{Proof.}}
	\setlength{\parindent}{0pt}
	}
	\tcolorbox[blanker,breakable,left=5mm,parbox=false,
    before upper={\parindent15pt},
    after skip=10pt,
	borderline west={1mm}{0pt}{clmscol!50!\colormode}]
}{
    \textcolor{clmscol!50!\colormode}{\hbox{}\nobreak\hfill$\blacksquare$}
    \endtcolorbox
}

\NewDocumentCommand{\clmp}{m+m+m}{
    \begin{claim*}{#1}{}
        #2
    \end{claim*}

    \begin{clmpf}
        #3
    \end{clmpf}
}


% ============================
% Proof
% ============================

\newenvironment{pf}[1]{%
  \def\pfcolor{#1}% Store the color in a macro
  \noindent{\it \textbf{Proof.}}%
  \tcolorbox[blanker,breakable,left=5mm,parbox=false,%
    before upper={\parindent15pt},%
    after skip=10pt,%
    borderline west={1mm}{0pt}{\pfcolor},%
    coltext=\TextColor
  ]%
  \setlength{\parindent}{0pt}%
}{%
  \textcolor{\pfcolor}{\hbox{}\nobreak\hfill$\blacksquare$}%
  \endtcolorbox%
}

\newtcolorbox{solution}{
  blanker,
  breakable,
  left=5mm,
  parbox=false,%
  before upper={\parindent15pt},%
  after skip=10pt,%
  borderline west={1mm}{0pt}{\pbcolormain},%
  coltext=\TextColor
}


% ============================
% Remark
% ============================


\NewDocumentCommand{\rmk}{+m}{
    {\it \color{rmkscol!100!white}#1}
}

\newenvironment{remark}{
    \par
    \vspace{5pt}
    \begin{minipage}{\textwidth}
        {\par\noindent{\textbf{Remark.}}}
        \tcolorbox[blanker,breakable,left=5mm,
        before skip=10pt,after skip=10pt,
        borderline west={1mm}{0pt}{rmkscol!80!\colormode},
        coltext=\TextColor
        ]
}{
        \endtcolorbox
    \end{minipage}
    \vspace{5pt}
}

\NewDocumentCommand{\rmkb}{+m}{
    \begin{remark}
        #1
    \end{remark}
}

\usepackage[all]{xypic}

% Define a new command for problems
\renewcommand{\headrule}{\color{\TextColor}\hrulefill}
\newcommand{\C}{\mathbb{C}}
\newcommand{\R}{\mathbb{R}}
\newcommand{\Q}{\mathbb{Q}}
\newcommand{\Z}{\mathbb{Z}}
\newcommand{\N}{\mathbb{N}}
\renewcommand{\P}{\mathbb{P}}
\newcommand{\F}{\mathbb{F}}
\newcommand{\contr}{\Rightarrow \Leftarrow}
\newcommand{\s}{\(\,\)}
\renewcommand{\t}[1]{\text{#1}}
\newcommand{\ang}[1]{\left\langle #1 \right\rangle}
\newcommand{\coker}{\mathrm{coker}}

% Meta data
\setstretch{1.2}

\begin{document}
\color{\TextColor}
\pagecolor{\colormode}
\setlength{\parindent}{0pt}
\pagestyle{fancy}
\fancyhf{}
\fancyhead[LOE]{\color{\TextColor}\slshape{\textbf{Vincent Ferrara}}}
\fancyhead[ROE]{\color{\TextColor}\thepage}
\begin{problem}{}
    Show that if $A$ is a Noetherian ring, then so is $A[[X]]$, the ring of formal power series with coefficients in $A$.
\end{problem}
\begin{problemproof}
    Assume $A$ is a Noetherian ring. 
Let $I$ be an ideal in $A[[X]]$.

Consider $I_n=\{a \in A \mid \exists f(X) \in I \t{ s.t. } f(X)=aX^n+a_{n+1}X^n+1+\dots\}$

Let $a,b \in I_n$ then $aX^n+\dots,bX^n+\dots \in I$ thus $(a+b)X^n+\dots \in I$ thus $(a+b) \in I_n$

Let $a \in I_n$ and $r \in A$ then $aX^n+\dots \in I$ thus $raX^n+\dots \in U$ thus $ra \in I_n$

Therefore, $I_n$ is an ideal of $A$ thus it is finitely generated.

Consider $I_n$ and $I_{n+1}$. If $a \in I_n$ then $aX^{n}+\dots \in I$ thus when you multiply by $X$ you get $aX^{n+1}+\dots \in I$ thus $a \in I_{n+1}$.

Therefore, $I_n \subseteq I_{n+1}$ thus we get an ascending chain of ideals. Since $A$ is Noetherian it has to stabilize at some point so let $I_N$ be the ideal in $A$ where it stabilizes (the lowest number). Thus, $I_N$ is finitely generated by $a_1,\dots,a_n$

Now pick $f_i(X)=a_iX^n+\dots$ 

Consider $(f_1,\dots,f_n)$

Now, take any \(f(X) \in I\)

If the lowest nonzero term of \(f(X)\) is at a degree less than \(N\), then since there are only finitely many such degrees (\(0, 1, \dots, N-1\)), we can generate all these parts by finitely many elements of \(I\) coming from the ideals \(I_0, I_1, \dots, I_{N-1}\). 

If the lowest nonzero term of \(f(X)\) is \(bX^k\) with \(k\ge N\), then \(b\in I_k = I_N\) (because the chain \(I_0\subseteq I_1\subseteq I_2\subseteq \dots\) stabilizes at \(I_N\)). Since \(I_N\) is finitely generated by \(a_1,\dots,a_n\), we have
\[
b = c_1a_1+\cdots+c_na_n.
\]
For each \(a_i\) there is a corresponding \(f_i(X) = a_iX^N+\dots \in I\). Multiplying \(f_i(X)\) by \(X^{k-N}\) yields a series with lowest nonzero term \(a_iX^k\), and then
\[
g(X)=c_1X^{k-N}f_1(X)+\cdots+c_nX^{k-N}f_n(X)
\]
has lowest nonzero term \(bX^k\) matching that of \(f(X)\). Subtracting \(g(X)\) from \(f(X)\) cancels the lowest nonzero term, and by repeating this process (which terminates in finitely many steps), we express \(f(X)\) as an \(A[[X]]\)-linear combination of finitely many generators. Thus, every \(f(X) \in I\) is generated by a finite set, proving that \(I\) is finitely generated and that \(A[[X]]\) is Noetherian.
\end{problemproof}

\begin{problem}{}
    Suppose that 
\[
\xymatrix{
0 \ar[r] & M' \ar[r]^{f_1}  \ar[d]^{\phi'}  & M \ar[r]^{g_1} \ar[d]^{\phi} & M'' \ar[r] \ar[d]^{\phi''} & 0 \\
0 \ar[r]& N' \ar[r]^{f_2} & N \ar[r]^{g_2} & N'' \ar[r]& 0 
}
\]
is a commutative diagram of $A$-modules with exact rows. Show that there is a long exact sequence
\[
\xymatrix{
0 \ar[r] & \ker(\phi')  \ar[r]  & \ker(\phi)  \ar[r]  & \ker(\phi'')  \ar[dll]_{\delta} \\
 & \coker(\phi')  \ar[r] & \coker(\phi)  \ar[r]& \coker(\phi'') \ar[r]& 0
}
\]
\rmkb{part of the problem is to define the maps in the long exact sequence; the only really interesting map is $\delta$. }
\end{problem}
\begin{problemproof}
    Define: $\varphi: \ker(\phi') \to \ker(\phi)$ s.t. $\varphi(x)=f_1(x)$

    Since $\phi(f_1(x))=f_2(\phi'(x))$, if $x \in \ker(\phi')$ then $\phi(f_1(x))=f_2(0)=0$ thus $f_1(x) \in \ker(\phi)$.

    Also since $f_1$ is injective $\varphi$ is injective thus $im(0)=\ker(\varphi)=\{0\}$\\

    Define: $\varphi': \ker(\phi) \to \ker(\phi'')$ s.t. $\varphi'(x)=g_1(x)$

    Since $\phi''(g_1(x))=g_2(\phi(x))$, if $x \in \ker(\phi)$ then $\phi''(g_1(x))=g_2(0)=0$ thus $g_1(x) \in \ker(\phi'')$

    Let $y \in im(\varphi)$ thus there is an $x \in \ker(\varphi')$ s.t. $\varphi(x)=f_1(x)=y$ thus $y \in im(f_1)=\ker(g_1)$. Therefore, $\varphi'(y)=g_1(y)=0$, $y \in \ker(\varphi')$

    Let $y \in \ker(\varphi')$ thus $\varphi'(y)=g_1(y)=0$ thus $y \in \ker(g_1)=im(f_1)$ thus there is an $x \in M'$ s.t. $f_1(x)=y$. Since $\phi(y)=f_2(\phi'(x))=0$ we know that $\phi'(x) \in \ker(f_2)=\{0\}$. Thus, $\phi'(x)=0$. Thus, $x \in \ker{\phi'}$. Therefore, $\varphi(x)=f_1(x)=y$. Thus, $x \in im(\varphi)$.

    Therefore $im(\varphi)=\ker(\varphi')$\\

    Let $x \in \ker(\phi'')$ since $g_1$ is surjective there is a non empty $A \subset M$ s.t. $g_1(A)=x$.

    Thus,
    \[g_2(\phi(A))=\phi''(g_1(A))=\phi''(x)=0\]
    Thus, $\phi(A) \subseteq \ker (g_2)=im(f_2)$

    Thus, $\exists B \subseteq N$ s.t. $f_2(B)=\phi(A)$
    
    Define: $\delta: \ker(\phi'') \to \t{coker}(\phi')$ s.t. $\delta(x)=B$ defined above.

    Well-Defined:

    Let $n_1,n_2 \in B$. Consider $f_2(n_1-n_2)=f_2(n_1)-f_2(n_2)$, since $f_2(B)=\phi(A)$ there is $m_1,m_2 \in A$ s.t. $\phi(m_1)=f_2(n_1)$ and $\phi(m_2)=f_2(n_2)$.

    Now consider $g_1(m_1-m_2)=x-x=0$. Thus, $m_1-m_2 \in \ker(g_1)=im(f_1)$ thus there is a $n \in N$ s.t.t $f_1(n)=m_1-n_2$.

    \[f_2(\phi'(n))=f_2(n_1-n_2) \implies \phi'(n)=n_1-n_2\]

    Thus, $n_1,n_2$ are in the same equivalence class thus $B$ is an is equivalence class.

    Let $y \in \ker(\delta)$. Since $g_1$ is surjective there is an $x \in M$ s.t. $g_1(x)=y$. Thus, by definition of $\delta$ since $0 \in im(\phi)$. We know that $0=f_2(0)=\phi(x)$ since $x \in g_1^{-1}(y)$. Thus, $x \in \ker(\phi)$. Thus, $\varphi'(x)=y$ thus $y \in im(\varphi')$.

    Let $y \in im(\varphi')$. Thus, there is an $x \in \ker(\phi)$ s.t. $\varphi'(x)=y$. Thus, $x \in g_1^{-1}(y)$, so look at $\delta(y)$. We know that $\delta(y)=B$ s.t. $f_2(b)=\phi(x)$ for $b \in B$.

    Since $\phi(x)=0$ and $f_2$ is injective we know that $b=0$ thus $b \in im(\phi)$. Thus, $B=im(\phi)$, since $B$ is an equivalence class. Therefore, $y \in \ker(\delta)$.

    Define: $\psi: \coker(\phi') \to \coker(\phi)$ s.t. $\psi([x])=[f_2(x)]$.

    Well-Defined:

    Let $[n_1]=[n_2]$ in $\coker(\phi')$

    Thus, $n_1-n_2 \in im(\phi')$, thus $\phi(m')=n_1-n_2$ for $m' \in M'$

    \[f_2(n_1)-f_2(n_2)=\phi(f_1(m'))\]

    Thus, $\psi([n_1])=\psi([n_2])$ since $f_2(n_1-n_2) \in im(\phi)$.

    Let $[n'] \in im(\delta)$ thus there is an $x \in \ker(\phi'')$ s.t. $\delta(x)=[n']$.

    By definition of delta we know that $f_2(n')=\phi(m)$ s.t. $g_1(m)=x$ for $m \in M$.

    Thus, $\psi(\delta(x))=[f_2(n')]=[\phi(m)]=0$.

    Let $[n'] \in \ker(\psi)$. Thus, this means that there is $m \in M$ s.t. $\phi(m)=f_2(n')$.

    Consider $x=g_1(m)$

    \[\phi''(x)=g_2(\phi(m))=g_2(f_2(n'))=0\]
    since $im(f_2)=\ker(g_2)$.

    Thus, $x \in \ker(\phi'')$ and we know that $f_2(n')=\phi(m) \implies \delta(x)=[n']$

    Thus, $im(\delta)=\ker(\psi)$.

    Define: $\psi': \coker(\phi) \to \coker(\phi'')$ s.t. $\psi([n])=[g_2(n)]$

    Well-defined as same proof as above.

    Let $[n] \in im(\psi)$ thus there is a $[n'] \in \coker(\psi)$ s.t. $\psi([n'])=[n]$. Thus $\psi'(\psi([n']))=[g_2(f_2(n'))]=0$ since $im(f_2)=\ker(g_2)$.

    Let $[n] \in \ker(\psi')$. Thus, $g_2(n)=\phi''(m'')$ for $m'' \in M''$.

    Since $g_1$ is surjective there is a $m \in M$ s.t. $g_1(m)=m''$.

    Consider:
    \[g_2(\phi(m))=\phi''(g_1(m))=g_2(n)\]

    Thus, $g_2(n-\phi(m))=0 \in \ker(g_2)=im(f_2)$ thus there is a $n' \in N'$ s.t.t $f_2(n')=n-\phi(m)$.

    Thus, $\psi([n'])=[n-\phi(m)]=[n]$ since $[\phi(m)]=0$.

    Thus $\ker(\psi')=im(\psi)$.

    Let $[n''] \in \coker(\phi'')$. Thus there is an $n \in N$ s.t. $g_2(n)=n''$. Thus $\phi'([n])=[n'']$. Thus, surjective.

    Thus, we have a long exact sequence.
\end{problemproof}

\addtocounter{subproblemcounter}{4}
\begin{problem}{}
    Show that $\mathcal{Z} (\mathcal{I} (Y)) = \overline{Y}$.
\end{problem}
\begin{problemproof}
    Since $\mathcal{Z}(\mathcal{I}(Y))=\{P \in A_k^n \mid f(P)=0 \t{ for all } f \in \mathcal{I}(Y)\}$.
    
    Thus, if $y \in Y$ and $f \in \mathcal{I}(Y)$ then $f(y)=0$ since $f(y)=0$ for all $y \in Y$. Thus, $y \in \mathcal{Z}(\mathcal{I}(Y))$. Thus, $Y \subseteq \mathcal{Z}(\mathcal{I}(Y))$.

    By definition of $\overline{Y} \implies \overline{Y} \subseteq \mathcal{Z}(\mathcal{I}(Y))$

    Let $C$ be a closed set containing $Y$. Thus, $C=\mathcal{Z}(T)$ for some $T \subseteq R$.

    Let $f \in T$ since $Y \subset C$ we know that for all $y \in Y$ satisfies $f(y)=0$ thus $f \in \mathcal{I}(Y)$. Thus, $T \subseteq \mathcal{I}(Y)$ by (a) $\mathcal{Z}(\mathcal{I}(Y)) \subseteq C$, thus $\mathcal{Z}(\mathcal{I}(Y))=\overline{Y}$ since it is a subset of every closed set containing $Y$.
\end{problemproof}

\begin{problem}{}
    {\bf The Nullstellensatz.} Continuation of the previous problem.
    The ideal
    \[
    \mathfrak{m}_P := (x_1 - \alpha_1, x_2 -\alpha_2, \ldots, x_n -\alpha_n)
    \]
    is maximal for any point $P=(\alpha_1, \ldots, \alpha_n) \in A^n_k$. 
    
    Here are two versions of the Nullstellensatz:

    (Weak version) Every maximal ideal in $R$ is of the form $\mathfrak{m}_P$ for some point $P \in A^n_k$. 

    (Strong version) $ \mathcal{I} (\mathcal{Z} (I)) = r(I)$, for any ideal $I\subseteq R$.
\end{problem}

\begin{subproblem}{}
    Show that the strong version implies the weak version.
\end{subproblem}
\begin{problemproof}
    Let $m$ be a maximal ideal in $R$ since $m \subseteq r(m) \implies m=r(m)$; therefore, $m=\mathcal{I}(\mathcal{Z}(m))$
    
    Assume for contradiction that $x,y \in \mathcal{Z}(m)$ s.t. $x \neq y$. Consider $\{x\}$. Since $\{x\} \subseteq \mathcal{Z}(m)$ by (a) we have that $\mathcal{I}(\mathcal{Z}(m))=m \subseteq \mathcal{I}(\{x\})$. This is a contradiction thus there is one point in $\mathcal{Z}(m)$. Call this point $P$
    \[m=\mathcal{I}(\mathcal(Z)(m))=\mathcal{I}(\{P\})=\{f \in R \mid f(P)=0\}=\mathfrak{m}_P\]
\end{problemproof}

\begin{subproblem}{}
    You may assume the strong version of the Nullstellensatz. Show that the assignments 
    \[
    Y \mapsto \mathcal{I} (Y), \quad I \mapsto \mathcal{Z}(I)
    \]
    give a bijection between closed subsets of $A^n_k$ and radical ideals of $R$. (An ideal $I$ in $R$ is called radical if $I=r(I)$.)
\end{subproblem}
\begin{problemproof}
    Let $I$ be a radical ideal of $R$ then consider $\mathcal{Z}(I)$ this is exactly a closed set of $A_k^n$
    
    Also, we have that $\mathcal{I}(\mathcal{Z}(I))=r(I)=I$
    
    Let $Y$ be a closed subset of $A_k^n$ then consider $\mathcal{I}(Y)$. Since $Y$ is closed we know that $Y=\mathcal{Z}(J)$ for some ideal $J \subseteq R$. Therefore, $\mathcal{I}(Y)=r(J)$. Consider $r(\mathcal{I}(Y))=r(r(J))=r(J)=\mathcal{I}(Y)$. Thus $\mathcal{I}(Y)$ is a radical ideal.
    
    Also, we have that $\mathcal{Z}(\mathcal{I}(Y))=\overline{Y}=Y$ 
    
    Therefore, these are inverses on closed subsets and radical ideals thus they form a bijection.
\end{problemproof}

\qh
\addtocounter{subproblemcounter}{3}
\begin{subproblem}{}
    Show that bijection in the previous problem restricts to a bijection between {\it irreducible} closed subsets of $A^n_k$ and {\it prime} ideals of $R$. 
\end{subproblem}
\begin{problemproof}{}
    Let $I$ be a prime ideal. Thus, $\mathcal{Z}(I)$ is closed. Assume for contradiction that $\mathcal{Z}(I)=\mathcal{Z}(J) \cup \mathcal{Z}(K)$ for $I \subsetneq J,K$ ideals of $R$.

    Thus, $\mathcal{Z}(I)=\mathcal{Z}(J) \cup \mathcal{Z}(K)=\mathcal{Z}(JK)$. Therefore, $r(I)=I=r(JK) \implies JK \subseteq r(JK)=I$.

    Since $I \subsetneq J,K$ pick
    \[ a\in J \setminus I \t{ and } b \in K \setminus I\]
    Thus, $a \in JK \subseteq I$ this is a contradiction of $I$ being prime. Thus, $\mathcal{I}$ is irreducible closed set.
    
    Also, we have that $\mathcal{I}(\mathcal{Z}(I))=r(I)=I$.

    Let $Y$ be an irreducible closed subset of $A_k^n$ then consider $\mathcal{I}(Y)$.

    Let $f,g \in \R$ s.t.
    \[(f \cdot g)(P)=0 \t{ for all } P \in Y\]
    thus $f \cdot g \in \mathcal{I}(Y)$

    Assume $f \notin \mathcal{I}(Y)$

    Consider $U=Y \setminus (Y \cap \mathcal{Z}(f))=\{P \in Y \mid f(P) \neq 0\}$. This is open in the subspace topology of $Y$ since $Y \cap \mathcal{Z}(f)$ is closed in $Y$.

    Let $P \in U \subseteq Y$. Since
    \[(f \cdot g)(P)=0\]
    This implies that $g(P)=0$ thus $U \subseteq \mathcal{Z}(g) \cap Y$
    
    $U$ is non-empty since $f \notin \mathcal{I}(Y)$ thus by (b) we know that $U$ is dense therefore, $U \subseteq \mathcal{Z}(g) \cap Y \implies \overline{U}=Y \subseteq \overline{\mathcal{Z}(g) \cap Y}=\mathcal{Z}(g) \cap Y \implies Y=\mathcal{Z}(g) \cap Y$.

    Therefore, $g \in \mathcal{I}(Y)$. Thus, $\mathcal{I}(Y)$ is prime.
    
    Also, we have that $\mathcal{Z}(\mathcal{I}(Y))=\overline{Y}=Y$.

    Thus, there is a bijection since these are inverses.
\end{problemproof}

\begin{subproblem}{}
    Let $I= (x_1 x_2) \subset k[x_1,x_2]$.  Is $I$ prime? Draw a picture showing what $\mathcal{Z} (I)$ looks like.  Is $\mathcal{Z} (I)$ irreducible? If not, write it as a union of two proper closed subsets. 
\end{subproblem}
\begin{problemproof}
    Consider $I=(x_1x_2)$. This is not prime since consider $x_1x_2 \in I$, but $x_1,x_2 \notin I$.
    
    Since if $x_1 \in I$ then there is an $f(x_1,x_2)$ s.t.
    \[x_1 = f(x_1,x_2)x_1x_2 \implies f(x_1,x_2)x_2=1\]
    This implies $x_2$ is a unit which it is not since deg$(x_2 \cdot f(x_1,x_2)) \geq 1 \neq \t{deg}(1)=0$.
    \begin{tikzpicture}[scale=1.5]
        % Draw the coordinate axes
        \draw[->, thick, blue] (-2.5,0) -- (2.5,0) node[right] {$x_1$};
        \draw[->, thick, red] (0,-2.5) -- (0,2.5) node[above] {$x_2$};

        % Draw the x1-axis (where x2 = 0) in blue
        \draw[thin, blue] (-2.5,0) -- (1.5,0) node[below right] {$x_2=0$};
        
        % Draw the x2-axis (where x1 = 0) in red
        \draw[thin, red] (0,-2.5) -- (0,2) node[above left] {$x_1=0$};
        
        % Mark the intersection point
        \fill (0,0) circle (2pt);
        \node[below right] at (0,0) {$(0,0)$};
        
        % Optional: Add a title below the picture
        \node at (0,-3) {$\mathcal{Z}(x_1x_2)=\{x_1=0\}\cup\{x_2=0\}=\mathcal{Z}(x_1) \cup \mathcal{Z}(x_2)$};
      \end{tikzpicture}
\end{problemproof}

\qh
\begin{subproblem}{}
    Artin, Ex. 11. 9.8.  Find all the ideals in the polynomial ring $\C [x,y] $ that contain $x^2 + y^2 -5$ and $xy -2$.
\end{subproblem}

Consider $\dfrac{\C[x,y]}{(x^2+y^2-5,xy-2)}$

Solving the system of equations we get these 4 points
\[(1,2), (2,1), (-1,-2), (-2,-1)\]
This is a closed subset of points since 

$Z((x^2+y^2-5,xy-2))=\{(1,2),(2,1),(-1,-2),(-2,-1)\}$

Therefore, we know that this is a radical ideal. Thus, it is the intersection of maximal ideals, since we know maximal ideals are of the form $(x-a,y-b)$ we know that:
\[(x^2+y^2-5,xy-2)=(x-1,y-2) \cap (x-2,y-1) \cap (x+1,y+2) \cap (x+2,y+1)\]

Since all maximal ideals are coprime by the chinese remainder theorems

\[\dfrac{\C[x,y]}{(x^2+y^2-5,xy-2)} \cong \dfrac{\C[x,y]}{(x-1,y-2)} \times \dfrac{\C[x,y]}{(x-2,y-1)} \times \dfrac{\C[x,y]}{(x+1,y+2)} \times \dfrac{\C[x,y]}{(x+2,y+1)}\]

\[\dfrac{\C[x,y]}{(x^2+y^2-5,xy-2)} \cong \C^4\]

Since just look at the map $\C[x,y] \to \C$ by sending $x,y$ to the roots of the bottom polynomials.

Thus, in this we know all the ideals are of the form some combination of $\C$ and $(0)$. Thus, we get $2^4=16$ ideals containing our original ideal.

These come from combinations of intersections of the maximal ideals we listed, so pick some of the maximal ideals and intersect them that gets us all 16.

\begin{subproblem}{}
    Artin, Ex. 11.9.12.  Prove in two ways that the polynomials
\begin{align*}
f_1 (x,y) &= y^2 + x^2 -2 \\
f_2 (x,y) &= yx-1 \\
f_3 (x,y) &= y^3 + 5yx^2 +1 
\end{align*}
generate the unit ideal in $\C[x,y]$: first,  by showing that they have no common zeros and second, by showing explicitly that 
you can write $1$ as $g_1 f_1 + g_2 f_2 + g_3 f_3$ for some polynomials $g_1, g_2, g_3$ in $\C[x,y]$. 
\end{subproblem}
\begin{problemproof}
Solve \(f_1(x,y)=0\) and \(f_2(x,y)=0\).

The equations
\[
\begin{cases}
x^2+y^2=2,\\[1mm]
xy=1,
\end{cases}
\]
are our starting point. Since \(xy=1\) we may (when \(x\neq0\)) solve for
\[
y=\frac{1}{x}.
\]
Substitute this into the first equation:
\[
x^2+\frac{1}{x^2}=2.
\]
Multiplying by \(x^2\) (which is nonzero) gives
\[
x^4+1=2x^2.
\]
Rearrange:
\[
x^4-2x^2+1=0,
\]
which factors as
\[
(x^2-1)^2=0.
\]
Thus, \(x^2=1\) so that \(x=\pm 1\) and consequently
\[
y=\frac{1}{x}=\pm 1.
\]
The two candidate common zeros (for \(f_1\) and \(f_2\)) are:
\[
(1,1)\quad\text{and}\quad (-1,-1).
\]

Test these candidates in \(f_3(x,y)=0\).

- For \((1,1)\):
  \[
  f_3(1,1)=1^3+5\cdot1^2\cdot1+1 = 1+5+1=7\neq 0.
  \]
- For \((-1,-1)\):
  \[
  f_3(-1,-1)=(-1)^3+5\cdot(-1)^2\cdot(-1)+1 = -1+5\cdot1\cdot(-1)+1=-1-5+1=-5\neq 0.
  \]

Since neither candidate makes \(f_3\) vanish, there is no point \((x,y)\in\C^2\) at which all three polynomials vanish simultaneously. By the weak Nullstellensatz it follows that
\[
(f_1,f_2,f_3)=\C[x,y],
\]

Consider
\begin{align*}
    f_3-yf_1&=4xy+2y+1\\
    4xy+2y+1-4f_2&=2y-3\\
    \intertext{Thus $2y-3 \in (f_1,f_2,f_3)$}
    2xy-3x-2f_1&=-3x+2\\
    \intertext{Thus $-3x+2 \in (f_1,f_2,f_3)$}
    (2y-3)(-3x+2)&=-6xy+4y+9x-6\\
    -6xy+4y+9x-6+6f_1&=4y+9x-12\\
    4y+9x-12+3(-3x+2)-2(2y-3)&=-12
    \intertext{Thus $-12 \in (f_1,f_2,f_3) \implies 1 \in (f_1,f_2,f_3)$}
\end{align*}
You can reconstruct the polynomials using the steps above and plugging in $f_1,f_2,f_3$ for all the polynomials used above.
\end{problemproof}
\end{document}